\documentclass[11pt,a4paper]{article}

\usepackage[margin=2.2cm]{geometry}
\usepackage{amsmath,amssymb,mathtools}
\usepackage{booktabs}
\usepackage{longtable}
\usepackage{array}
\usepackage{xcolor}
\usepackage{siunitx}
\usepackage{microtype}
\usepackage{enumitem}
\usepackage[hidelinks]{hyperref}

\newcommand{\Vm}{V_{m}}
\newcommand{\Iion}{I_{\text{ion}}}
\newcommand{\tact}{t_{\text{act}}^{P8}}
\newcommand{\code}[1]{\texttt{\small #1}}

\title{Cross-solver comparative summary\\
       \large Manufactured verification and Niederer electro-activation
       benchmark across the five high-order solvers
       in \texttt{cardiacFoam-pc}}
\author{Pablo}
\date{\today}

\begin{document}
\maketitle

\begin{abstract}
This document is an interim cross-report synthesis of the five solver
reports currently maintained under
\code{applications/solvers/}: two method-of-manufactured-solutions
(MMS) reports that verify the discretisation, and three
electro-activation reports that exercise the same discretisation on
the Niederer-2012 monodomain benchmark. It is not a substitute for the
individual reports; it is meant to answer, in one place, the six
practical questions that come up when looking at the runs as a whole:
does the method work, what does the benchmark converge to, which
discretisation combinations are most cost-effective for that value,
what those combinations cost, which combination is the most accurate
and which is the most efficient, and what is still missing from the
present data. All numbers below are transcribed from the per-solver
result tables that exist at the time of writing; no new runs were
performed for this summary.
\end{abstract}

\tableofcontents
\bigskip

% ================================================================== %
\section{Scope and methods catalogue}
\label{sec:scope}
% ================================================================== %

The five reports cover the same discretisation toolbox --- standard
cell-centred finite-volume (FV) diffusion or a Local
Regression-Estimator (LRE) high-order reconstruction; cell-centred or
LRE-quadrature ionic source; lumped or consistent LRE mass matrix;
$\theta$-scheme time integration; and one of three nonlinear coupling
strategies (\code{Picard}, \code{diagonalIion}, \code{JFNK}) when the
PDE is implicit --- applied to two problems:

\begin{itemize}[leftmargin=*]
\item A manufactured PDE--ODE system with a known exact solution
($\Vm^\star$, $u_1^\star$, $u_2^\star$, $u_3^\star$) and a smooth
nonlinear source $\Iion^{\text{MMS}}(\Vm,\mathbf{u})$.
\item The Niederer-2012 monodomain benchmark with the TNNP ionic
model, with $\tact$ as the diagnostic scalar.
\end{itemize}

\begin{center}\footnotesize
\begin{longtable}{p{0.27\linewidth}p{0.10\linewidth}p{0.18\linewidth}p{0.16\linewidth}p{0.20\linewidth}}
\caption{Catalogue of the five solver reports.}\\
\toprule
Solver & Problem & Time scheme & Nonlinear & Dimensions covered \\
\midrule\endfirsthead\toprule
Solver & Problem & Time scheme & Nonlinear & Dimensions covered \\
\midrule\endhead
\code{highOrderManufacturedFDAExplicit}     & MMS  & forward Euler (explicit) & ---                                        & 2D ($N=10\!\dots\!100$), 3D ($N=10\!\dots\!30$) \\
\code{highOrderManufacturedFDAImplicitJfNK} & MMS  & bE / cN ($\theta$-scheme)& Picard / diagonalIion / JFNK               & 2D ($N=10\!\dots\!100$), 3D ($N=10\!\dots\!30$) \\
\code{highOrderElectroActivationFoamExplicit}     & Niederer / TNNP & forward Euler & ---                       & 2D ($\Delta x\!\in\![0.5,0.05]$ mm), 3D ($\Delta x\!\in\![0.5,0.1]$ mm) \\
\code{highOrderElectroActivationFoamImplicit}     & Niederer / TNNP & cN, $\theta=\tfrac{1}{2}$ & single-sweep Picard ($k=1$) & 2D ($\Delta x\!\in\![0.5,0.05]$ mm, $\Delta t\!\in\![10^{-3},10^{-6}]$ s), 3D ($\Delta x\!\in\![0.5,0.1]$ mm) \\
\code{highOrderElectroActivationFoamImplicitJfNK} & Niederer / TNNP & cN, $\theta=\tfrac{1}{2}$ & Picard / diagonalIion / JFNK & 2D ($\Delta x\!\in\![0.5,0.1]$ mm, $\Delta t = 10^{-4}$ s), 3D ($\Delta x = 0.5$ mm) \\
\bottomrule
\end{longtable}\end{center}

Throughout this document an LRE pair is written
$(\code{p}_V, \code{p}_I)$, where the first slot
($\code{p}_V$, for $\Vm$) is the high-order level used for the
diffusion operator (and, when consistent mass is on, for the mass
matrix) and the second slot ($\code{p}_I$, for $\Iion$) is the level
used for the ionic source and ODE quadrature. \code{NO} means the standard FV / cell-centre operator.
\code{NC} marks any cell where the corresponding run has not yet been
done.

% ================================================================== %
\section{Convergence verification (MMS)}
\label{sec:mms}
% ================================================================== %

\subsection{What the LRE level is supposed to deliver}

The two MMS reports state the same theoretical-accuracy ladder:
\code{NO} and \code{p1} target second-order spatial accuracy,
\code{p2} third-order, and \code{p3} fourth-order. The cell-centred FV
operator itself is second-order, so the cell-centred error column can
only saturate at $\sim 2$; the formal LRE order can only be seen in
the \emph{reconstructed} (Gauss-quadrature) error columns. The fitted
order is also the minimum of the spatial order and the temporal /
nonlinear-tolerance order; an aggressive LRE polynomial does not
recover its formal order unless $\Delta t$, the mass matrix, and the
nonlinear residual all permit it.

\subsection{Explicit MMS --- observed orders}

From
\code{highOrderManufacturedFDAExplicit/report/tables/explicit\_orders\_2D.tex}
and \code{explicit\_orders\_3D.tex}, fitted on $N=10\dots100$ (2D) and
$N=10\dots30$ (3D), at the finest $\Delta t = 10^{-6}$~s:

\begin{center}\footnotesize
\begin{longtable}{lcccccc}
\caption{Explicit MMS: observed convergence orders at
$\Delta t = 10^{-6}$~s. Cell / G2 = cell-centred / reconstructed
Gauss-quadrature error. Targets are written as
$\text{Vm}/\Iion$ (LRE polynomial order).}\\
\toprule
$(\code{p}_V,\code{p}_I)$ & target & $\Vm$ cell (2D) & $\Vm^G$ (2D) & $\Vm$ cell (3D) & $\Vm^G$ (3D) \\
\midrule\endfirsthead\toprule
$(\code{p}_V,\code{p}_I)$ & target & $\Vm$ cell (2D) & $\Vm^G$ (2D) & $\Vm$ cell (3D) & $\Vm^G$ (3D) \\
\midrule\endhead
(\code{NO},\code{NO}) & 2/2 & 1.997 & 1.997 & 1.979 & 1.979 \\
(\code{p1},\code{p1}) & 2/2 & 1.912 & 2.147 & 2.008 & 2.063 \\
(\code{p2},\code{p2}) & 3/3 & 2.244 & 2.762 & 2.021 & 3.289 \\
(\code{p3},\code{p3}) & 4/4 & 2.369 & 3.091 & 3.108 & 3.869 \\
\bottomrule
\end{longtable}\end{center}

Reading the table: \code{NO}/\code{NO} and \code{p1}/\code{p1} are
clean second-order in both columns; \code{p2}/\code{p2} reaches its
target third order in the reconstructed column in 3D (3.289) and is
close in 2D (2.762); \code{p3}/\code{p3} hits 3.869 in 3D --- formal
fourth order is approached but not fully recovered with explicit
forward-Euler at $\Delta t = 10^{-6}$~s. The cell-centred Vm column
saturates at $\sim 2$ in 2D as expected.

\subsection{Implicit MMS --- where formal fourth order is recovered}

The implicit JFNK MMS report sweeps the same LRE pairs with both time
schemes (\code{bE}/\code{cN}), both mass matrices
(\code{lumped}/\code{consistent}), and all three nonlinear methods.
The cleanest case for fourth order is \code{p3}/\code{p2} with
\code{crankNicolson}, consistent mass, and $\Delta t = 10^{-3}$~s
(\code{tables/results\_2D\_*\_dt\_1p0em03.tex}, line~65):

\begin{center}\footnotesize
\begin{longtable}{p{0.20\linewidth}cccccc}
\caption{Implicit MMS 2D, \code{p3}/\code{p2}, $\Delta t = 10^{-3}$~s.
Orders fitted over $N=10\dots100$. Cost is the
$t_\Sigma$ / RSS$_\Sigma$ row sum from the same table.}\\
\toprule
Configuration & $\Vm$ & $\Vm^G$ & $u_1$ & $u_1^G$ & $u_2$ & $u_2^G$ \\
\midrule\endfirsthead\toprule
Configuration & $\Vm$ & $\Vm^G$ & $u_1$ & $u_1^G$ & $u_2$ & $u_2^G$ \\
\midrule\endhead
\code{bE} / consistent     & 2.434 & 3.214 & 2.000 & 3.534 & 2.000 & 3.532 \\
\code{cN} / consistent     & 3.539 & 3.993 & 1.999 & 4.003 & 1.997 & 4.010 \\
\code{cN} / lumped         & 2.369 & 3.091 & 1.999 & 3.441 & 1.996 & 3.435 \\
\bottomrule
\end{longtable}\end{center}

Three concrete observations from those three rows:

\begin{itemize}[leftmargin=*]
\item \textbf{Backward Euler caps to first order in time}, which here
collapses the reconstructed Vm fit from $3.993 \rightarrow 3.214$; the
state-reconstructed columns survive because their dominant error
component is spatial, not temporal.
\item \textbf{Lumped mass caps the reconstructed Vm order}
($3.993 \rightarrow 3.091$); the consistent LRE mass matrix is what
lets the reconstruction respect its own quadrature.
\item \textbf{Crank--Nicolson with consistent mass is the only
combination that recovers the formal fourth-order target} for both
$\Vm^G$ and the reconstructed state errors ($3.993, 4.003, 4.010$).
\end{itemize}

Comparing the three nonlinear methods on the very same
\code{p3}/\code{p2}, \code{cN}, consistent, $\Delta t = 10^{-3}$~s row:

\begin{center}\footnotesize
\begin{tabular}{lcccc}
\toprule
Nonlinear method  & $\Vm$ & $\Vm^G$ & $t_\Sigma$ [s] & RSS$_\Sigma$ [MB] \\
\midrule
\code{Picard}       & 3.538 & 3.992 & 1730.1 & 19265.3 \\
\code{diagonalIion} & 3.538 & 3.992 & 1661.8 & 19310.4 \\
\code{JFNK}         & 3.539 & 3.993 &  798.3 & 18963.3 \\
\bottomrule
\end{tabular}\end{center}

The three methods produce identical orders to three decimals --- as
expected, since the MMS source is smooth and the nonlinear coupling
error is small. They differ purely in cost: JFNK is $\sim 2.2\times$
cheaper than Picard / diagonalIion in this MMS test. \emph{This is
the opposite of the electro-activation picture below}, and the
difference is explained in \S\ref{sec:efficient}.

\subsection{Bottom line --- does the method work?}

Yes. The MMS evidence is consistent: the formal LRE polynomial order
is recovered whenever the rest of the discretisation does not block
it (Crank--Nicolson, consistent mass, $\Delta t$ small enough,
nonlinear residual converged). The MMS does not stress-test the
non-smoothness of an activation front or the stiffness of the TNNP
ionic model; those stress tests are the role of the Niederer reports.

% ================================================================== %
\section{Electro-activation benchmark (Niederer)}
\label{sec:niederer}
% ================================================================== %

This is the same monodomain PDE + TNNP-2004 ionic model in three
geometries: a 2D code-development reduction ($20\times 8$ mm, single
cell in $z$, planar fibres; \emph{not} a published Niederer geometry)
and the published 3D Niederer slab ($20\times 3\times 7$ mm,
anisotropic fibre rotation, $40\times 6\times 14$ cells at
$\Delta x = 0.5$~mm). The scalar diagnostic is the P8 activation time
$\tact$~[ms]: the time at which a designated benchmark point at the
far end of the slab first crosses $\Vm = V_{\text{thr}}$.

\subsection{2D internal-consistency floor}

Neither the 2D explicit nor the 2D implicit report can compare $\tact$
against an external reference (the 2D reduction is not in the
published benchmark literature). Instead each report reports an
\emph{internal-consistency floor}: the band in which the high-order
LRE pairs all agree at the finest mesh / smallest $\Delta t$.

\begin{center}\footnotesize
\begin{longtable}{p{0.30\linewidth}p{0.22\linewidth}p{0.22\linewidth}p{0.18\linewidth}}
\caption{2D internal-consistency floor for $\tact$ across the three
electro-activation reports.}\\
\toprule
Solver & Floor for $\tact$ [ms] & Where it is reached & Reached by \\
\midrule\endfirsthead\toprule
Solver & Floor for $\tact$ [ms] & Where it is reached & Reached by \\
\midrule\endhead
\code{...FoamExplicit}        & $[44.5,\,46.0]$  & $\Delta x \le 0.1$~mm        & \code{(NO,p2)}, \code{(NO,p3)}, \code{(p1,p2)}, \code{(p1,p3)}, \code{(p3,p3)} \\
\code{...FoamImplicit}        & $[45.04,\,45.80]$& $\Delta x = 0.1$~mm, $\Delta t = 10^{-6}$~s & every HO pair we ran at the finest column \\
\code{...FoamImplicitJfNK}    & (no floor yet)   & 2D finest mesh available is $\Delta x = 0.1$~mm and only for \code{Picard} & \code{(NO,NO)}=$42.76$, \code{(NO,p1)}=$41.84$; HO pairs were not run at $\Delta x = 0.1$~mm yet \\
\bottomrule
\end{longtable}\end{center}

The first two solvers agree to within $\sim 0.5$~ms on a floor of
\textbf{$\tact \approx 45$--$46$~ms} for this 2D reduction. The JFNK
report has not yet been pushed deep enough into the high-order /
fine-mesh corner to corroborate that floor; the only $\Delta x = 0.1$~mm
entries it has are low-order Picard rows that sit slightly below it
(\code{(NO,NO)} drops to $42.76$~ms when $\Delta t$ goes from
$10^{-5}$ to $10^{-4}$~s, which is the single-Picard linearisation
effect discussed in \S\ref{sec:efficient}).

\subsection{3D current best estimate}

The 3D data are much sparser. Tabulating what each electro-activation
report has at $\Delta x = 0.5$~mm (the only mesh that exists in all
three):

\begin{center}\footnotesize
\begin{longtable}{p{0.10\linewidth}p{0.20\linewidth}p{0.18\linewidth}p{0.18\linewidth}p{0.18\linewidth}}
\caption{3D $\tact$ [ms] at $\Delta x = 0.5$~mm across the three
electro-activation reports. \code{NC} marks LRE pairs that were not
yet run with that solver.}\\
\toprule
Pair & Explicit ($\Delta t_{\text{req}}=5\!\cdot\!10^{-5}$~s)
     & Implicit ($\Delta t=5\!\cdot\!10^{-5}$~s)
     & JfNK-Picard ($\Delta t=10^{-4}$~s)
     & JfNK-diagIion ($\Delta t=10^{-4}$~s) \\
\midrule\endfirsthead\toprule
Pair & Explicit ($\Delta t_{\text{req}}=5\!\cdot\!10^{-5}$~s)
     & Implicit ($\Delta t=5\!\cdot\!10^{-5}$~s)
     & JfNK-Picard ($\Delta t=10^{-4}$~s)
     & JfNK-diagIion ($\Delta t=10^{-4}$~s) \\
\midrule\endhead
\code{(NO,NO)} & 141.5 & 142.7 & 136.7 & 136.6 \\
\code{(NO,p1)} &  36.4 &  38.0 &  33.5 &  33.4 \\
\code{(NO,p3)} & NC    & NC    &  30.6 &  30.5 \\
\code{(p3,NO)} & NC    & NC    & 149.1 & 149.1 \\
\code{(p3,p1)} & NC    & NC    &  32.4 &  32.4 \\
\code{(p3,p3)} &  32.6 &  34.5 &  29.6 & NC    \\
\bottomrule
\end{longtable}\end{center}

The qualitative pattern is the same in all three reports: turning on
LRE on the ionic source alone collapses the activation time from
$\sim 140$~ms (low-order) to $\sim 30$--$38$~ms; the matched HO pair
adds another $\sim 2$--$5$~ms drop. The three solvers agree within
$\sim 0$--$8$~ms across all rows, with the dominant offset coming
from $\Delta t$ ($5\!\cdot\!10^{-5}$ vs $10^{-4}$~s) and from the
single-Picard linearisation in the non-JFNK implicit branch.

\textbf{The 3D problem does not yet have enough mesh refinement to
claim a converged $\tact$.} The high-order pairs cluster in
$[29.6,\,34.5]$~ms at $\Delta x = 0.5$~mm, but the only finer 3D mesh
that exists is for the low-order pair \code{(NO,NO)} (implicit:
$55.6$~ms at $0.2$~mm, $47.3$~ms at $0.1$~mm). Until at least one HO
pair is rerun at $\Delta x \le 0.25$~mm, the 3D internal-consistency
floor remains open.

\subsection{Cross-method agreement on the same configuration}

The implicit-JFNK report contains the only side-by-side comparison of
\code{Picard}, \code{diagonalIion}, and \code{JFNK} on the same
electro-activation case. At $\Delta x = 0.5$~mm, $\Delta t = 10^{-4}$~s:

\begin{center}\footnotesize
\begin{tabular}{lccc}
\toprule
Pair & Picard [ms] & diagonalIion [ms] & JFNK [ms] \\
\midrule
\code{(NO,NO)} & 153.32 & 153.16 & 153.13 \\
\code{(NO,p3)} &  32.46 &  32.37 &  32.36 \\
\code{(p3,p3)} &  33.09 &  33.00 &  33.00 \\
\bottomrule
\end{tabular}\end{center}

The three nonlinear methods agree to within $0.1$--$0.2$~ms on every
row available. In other words, on the Niederer benchmark the
\emph{choice of nonlinear method is essentially accuracy-free} at
$\Delta t = 10^{-4}$~s; it is a cost decision (see \S\ref{sec:cost}).

% ================================================================== %
\section{Most cost-effective combinations}
\label{sec:efficient}
% ================================================================== %

\subsection{High-order on the source alone is the dominant accuracy lever}

The single largest accuracy gain per unit of extra cost comes from
turning on LRE quadrature for the ionic source while leaving the
diffusion at the standard FV operator: the \code{(NO,p2)} or
\code{(NO,p3)} pair. Concrete example from the implicit 2D table
(\code{...FoamImplicit/report/results\_tables.tex}, $\Delta t = 10^{-5}$~s):

\begin{center}\footnotesize
\begin{tabular}{lccccc}
\toprule
Pair & $\Delta x = 0.5$ & $0.25$ & $0.125$ & $0.1$ & $0.05$ [mm] \\
\midrule
\code{(NO,NO)} $\tact$ [ms]              & 157.69 & 67.69 & 50.65 & 48.55 & 46.10 \\
\code{(NO,p2)} $\tact$ [ms]              &  35.26 & 43.20 & 45.15 & 45.42 & --- \\
\code{(NO,p2)/(NO,NO)} wall-clock ratio  & 2.0    & 3.0   & 5.5   & 4.5   & --- \\
\bottomrule
\end{tabular}\end{center}

Reading row 1: the low-order pair only enters the floor band at the
two finest meshes. Reading row 2: \code{(NO,p2)} reaches a value
within $\sim 2$~ms of the floor already at $\Delta x = 0.5$~mm and is
inside the band at $\Delta x = 0.125$~mm. The cost ratio (row 3) is
between $\sim 2\times$ and $\sim 5\times$ for the same target
accuracy --- in other words, far cheaper than the order-of-magnitude
mesh refinement that the \code{(NO,NO)} branch would need to match it.

The same pattern holds in the explicit report (\code{(NO,p2)} reaches
$34.7 \rightarrow 44.95$~ms over the mesh sweep) and in the
implicit-JFNK report (\code{(NO,p3)} reaches $32.4 \rightarrow$
finer-mesh values within the band).

\subsection{LRE on the diffusion only is the worst-of-both-worlds}

The \code{(p$k$, NO)} branch --- high-order LRE Laplacian on top of a
piecewise-constant ionic source --- is documented in the explicit
report as numerically pathological on coarse meshes. The Niederer
explicit 2D table shows \code{(p1,NO)} at $\Delta x = 0.5$~mm giving
$\tact = 14.25$~ms (a spurious very-early crossing) and
\code{(p2,NO)} giving $0.00$~ms (the front never crossed in the
simulated window). Same pattern in the implicit-JFNK 2D table:
\code{(p3,NO)} at $\Delta x = 0.5$~mm gives $180.4$--$180.5$~ms across
all three nonlinear methods, $\sim 20$\% \emph{worse} than the
\code{(NO,NO)} baseline at the same mesh. This branch should be
treated as a diagnostic, not a production setting.

\subsection{The best combination versus the most efficient combination}

Putting the two together:

\begin{itemize}[leftmargin=*]
\item \textbf{Best combination} (highest fidelity reference). A matched
high-order pair (\code{p3}/\code{p3} or \code{p3}/\code{p2}) with
Crank--Nicolson, consistent mass, the finest mesh in the table, and
$\Delta t = 10^{-6}$~s on the implicit solver. This is the
combination that \emph{defines} the internal-consistency floor in 2D
($45.0$--$45.8$~ms). The single-Picard cost in this corner is
$\sim 700$--$900$~min wall-clock on the finest 2D mesh per case
(\code{...FoamImplicit/results\_tables.tex} wall-clock block).
\item \textbf{Most efficient combination} (highest accuracy per unit
cost). \code{(NO,p2)} or \code{(NO,p3)} on a coarse mesh, with
\code{Picard} as the nonlinear method when the implicit branch is
needed --- or with explicit forward-Euler when the CFL clamp at
$\Delta x = 0.5$~mm ($\Delta t_{\text{CFL}} \approx 8.7\!\cdot\!10^{-5}$~s
in 3D) is acceptable. Picard and diagonalIion agree on $\tact$ in
both 2D and 3D, and Picard is typically $\sim 10$--$20$\% faster
than diagonalIion at this accuracy level on the Niederer problem.
\end{itemize}

JFNK on the electro-activation problem is \emph{not} the most
efficient choice, despite being so on the MMS problem. The reason is
that the dominant cost of a JFNK iteration on the electro-activation
problem is the per-Krylov-vector TNNP ODE re-solve (RKF45, 17 states,
adaptive sub-stepping); each GMRES matvec re-integrates the ionic ODE
over $[t^n, t^{n+1}]$ on every cell or quadrature point, which is
much more expensive than the cheap manufactured source of the MMS
test. Concrete penalty from the implicit-JFNK 2D wall-clock table
(\code{...FoamImplicitJfNK/.../solver\_report.tex} Table~6) at
$\Delta x = 0.125$~mm, \code{(NO,p3)}: Picard
$402$~min vs JFNK $3407$~min (\textbf{$\sim 8.5\times$} penalty;
\code{diagonalIion} sits between them at $533$~min).

% ================================================================== %
\section{Computational cost cross-table}
\label{sec:cost}
% ================================================================== %

The following table lists four landmark configurations, one per row,
and reports their wall-clock cost (in minutes) and peak resident set
size (in MiB) on the cheapest available solver / nonlinear method.
The numbers come directly from the per-solver result tables.

\begin{center}\footnotesize
\begin{longtable}{p{0.34\linewidth}p{0.18\linewidth}p{0.12\linewidth}p{0.10\linewidth}p{0.16\linewidth}}
\caption{Cost cross-table for four landmark configurations. Time is
elapsed wall-clock in minutes (rounded). \code{NA} = configuration
not applicable to that solver.}\\
\toprule
Configuration & Solver / method & 2D time & 2D RSS & 3D time \\
\midrule\endfirsthead\toprule
Configuration & Solver / method & 2D time & 2D RSS & 3D time \\
\midrule\endhead
Cheapest baseline:           & Explicit                 & 1 min        & 148 MiB  & 3 min                   \\
\code{(NO,NO)}, $\Delta x = 0.5$~mm
                             & Implicit / Picard, $\Delta t = 10^{-3}$ & 1 min        & 148 MiB  & 4 min ($\Delta t = 5e\!-\!5$)  \\
                             & JfNK / Picard, $\Delta t = 10^{-4}$     & 15 min       & 148 MiB  & 80 min                 \\
\midrule
Cheapest accurate:           & Explicit (CFL-clamped)   & 1 min        & 152 MiB  & 11 min ($\Delta x = 0.5$, $\Delta t_{\text{CFL}}$) \\
\code{(NO,p2)} or \code{(NO,p3)}, $\Delta x = 0.5$~mm
                             & Implicit / Picard, $\Delta t = 10^{-5}$ & 2--3 min     & 155 MiB  & 14 min ($\Delta x = 0.5$, $\Delta t = 5e\!-\!5$)             \\
                             & JfNK / Picard, $\Delta t = 10^{-4}$     & 23 min       & 159 MiB  & 168 min                \\
\midrule
Mid-range converged:         & Explicit                 & 42 min       & 234 MiB  & NA                     \\
\code{(NO,p2)}, $\Delta x = 0.125$~mm
                             & Implicit / Picard, $\Delta t = 10^{-5}$ & 44 min       & 269 MiB  & NA                     \\
                             & JfNK / Picard, $\Delta t = 10^{-4}$     & 402 min      & 288 MiB  & NA                     \\
\midrule
Reference run (defines floor): & Explicit ($\Delta x = 0.05$~mm) & 1167 min  & 961 MiB  & NA \\
\code{(p2,p2)}, finest mesh
                             & Implicit / Picard, $\Delta x = 0.05$, $\Delta t = 10^{-3}$ & 246 min & 1429 MiB & NA \\
                             & Implicit / Picard, $\Delta x = 0.1$, $\Delta t = 10^{-6}$  & 730 min & 466 MiB  & NA \\
\bottomrule
\end{longtable}\end{center}

The two general lessons that fall out of this table are:

\begin{itemize}[leftmargin=*]
\item \textbf{Crossing into a finer mesh costs $10$--$30\times$ more
than upgrading the LRE pair on a coarse mesh}; e.g. going from
\code{(NO,p2)} at $\Delta x = 0.5$~mm (2~min) to \code{(NO,p2)} at
$\Delta x = 0.125$~mm (44~min) is a $\sim 20\times$ cost step, and
buys only $\sim 10$~ms of accuracy ($\tact$ moves from $35.3$ to
$45.2$~ms, i.e. from below to inside the floor band).
\item \textbf{The JFNK branch on the electro-activation problem
multiplies wall-clock cost by $\sim 5$--$8\times$ vs Picard at the
same accuracy}; see for instance the
\code{(NO,p3)}/$\Delta x = 0.125$~mm row of the JFNK 2D table cited
in \S\ref{sec:efficient}. Memory cost is essentially method-independent.
\end{itemize}

% ================================================================== %
\section{Gaps and pending runs}
\label{sec:gaps}
% ================================================================== %

The largest \code{NC} blocks in the current data, listed in the order
in which closing them most increases the value of the reports:

\begin{enumerate}[leftmargin=*]
\item \textbf{JFNK 3D sweep has not been run at all.} The
implicit-JFNK 3D table has \code{Picard} and \code{diagonalIion}
columns but no \code{JFNK} column. Filling at least the six 2D-row
configurations at $\Delta x = 0.5$~mm closes the cross-method
comparison in 3D.
\item \textbf{One finer 3D mesh on a high-order pair.} The implicit
3D table only has \code{(NO,NO)} at $\Delta x \in \{0.5, 0.2, 0.1\}$~mm;
no high-order pair has been run below $\Delta x = 0.5$~mm in 3D. A
single run of \code{(NO,p2)} or \code{(p3,p3)} at
$\Delta x = 0.25$~mm or $0.2$~mm (cheapest on explicit forward-Euler
at $\Delta t = 5\!\cdot\!10^{-6}$~s, or on Picard at
$\Delta t = 10^{-4}$~s) would let us define the 3D
internal-consistency floor.
\item \textbf{Missing \code{(p3,p3)} diagonalIion 3D row} in the
implicit-JFNK 3D table.
\item \textbf{Finest-mesh entries for \code{(NO,p2)} and \code{(NO,p3)}
in the explicit 2D table} at $\Delta x = 0.05$~mm are \code{NC};
these are the runs that would let us claim the explicit 2D floor
\emph{also} from the most cost-effective LRE pairs, not only from the
matched high-order pairs.
\item \textbf{Several Picard-only entries at $\Delta t = 10^{-6}$~s,
$\Delta x = 0.05$~mm} in the implicit (non-JFNK) 2D table are
\code{NC} (mostly the \code{p3} rows). These are the heaviest cases
in the entire 2D study; they are not blocking conclusions, only
documentation.
\item \textbf{JFNK 2D fine-mesh rows.} JFNK 2D currently stops at
$\Delta x = 0.125$~mm; only \code{Picard} reaches $\Delta x = 0.1$~mm.
Two or three JFNK runs at the same mesh as the existing Picard rows
would close the JFNK 2D internal-consistency picture.
\end{enumerate}

Items (1) and (2) are the two runs that most change the qualitative
conclusion of this summary; everything else is documentation.

% ================================================================== %
\section{Bottom-line conclusions}
\label{sec:conclusions}
% ================================================================== %

Answering, in order, the six questions that prompted this document:

\begin{enumerate}[leftmargin=*]
\item \textbf{Do the methods work?} Yes. The MMS reports recover the
formal LRE polynomial order on every clean combination
(Crank--Nicolson, consistent mass, $\Delta t$ small enough). The three
Niederer-electro-activation reports agree with each other in the
overlap regions; cross-method differences are explained by their
respective discretisation choices, not by a bug.
\item \textbf{What does the electro-activation problem converge to?}
\textbf{2D:} $\tact \approx 45$--$46$~ms (internal-consistency floor
shared by the explicit and the single-Picard implicit reports). The
JFNK report has not yet run deep enough to corroborate this floor
from a different time-step / nonlinear path. \textbf{3D:} the answer
is currently \emph{open}. The high-order pairs cluster in
$29.6$--$34.5$~ms at $\Delta x = 0.5$~mm but the 3D mesh sweep is too
shallow to define a converged value.
\item \textbf{Which combinations are the most cost-effective for that
value?} \code{(NO,p2)} or \code{(NO,p3)} --- LRE quadrature on the
ionic source alone, with the standard FV diffusion. They reach the
internal-consistency band already at $\Delta x = 0.125$~mm and cost
$\sim 5\times$ less than the matched high-order pairs at the same
mesh.
\item \textbf{What do those combinations cost?} On the implicit
(single-Picard) 2D table at $\Delta t = 10^{-5}$~s: $\sim 2$~min at
$\Delta x = 0.5$~mm, $\sim 44$~min at $\Delta x = 0.125$~mm,
$\sim 243$~min at $\Delta x = 0.05$~mm for \code{(NO,p2)}; peak RSS
stays under $1$~GiB. JFNK multiplies the wall-clock by
$\sim 5$--$8\times$ for the same configuration on the
electro-activation problem.
\item \textbf{Which is the best combination?} The matched
high-order pair \code{p3}/\code{p3} (or \code{p3}/\code{p2}) with
\code{cN} + consistent mass + finest mesh + smallest $\Delta t$. This
is the combination that \emph{defines} the floor; it is also the
single most expensive cell of the entire study.
\item \textbf{Which is the most efficient combination?}
\code{(NO,p2)} or \code{(NO,p3)} on a coarse mesh, with \code{Picard}
when the implicit branch is needed. On the electro-activation problem
the JFNK branch is \emph{not} efficient (the per-Krylov ODE re-solve
of the TNNP cell model dominates); this is the qualitative reverse of
the MMS picture in \S\ref{sec:mms}, where JFNK is the cheapest of
the three.
\end{enumerate}

\paragraph{A note on Crank--Nicolson order.}
The implicit (non-JFNK) electro-activation report documents that with
a \emph{single} Picard sweep per time step, the formal
$\theta = \tfrac{1}{2}$ second-order temporal behaviour is not
recovered (\code{$\tact$} on the finest mesh halves between
$\Delta t = 10^{-3}$ and $10^{-4}$~s, rather than quartering as the
asymptote would predict). This is a linearisation-error signature,
not a Crank--Nicolson failure; the implicit-JFNK solver fixes it by
genuinely solving the coupled nonlinear residual. The implicit
report's stated next step --- enabling a tolerance-based stop on the
Picard residual with $k \ge 2$ sweeps --- is the cheapest way to
recover CN second order without paying the JFNK ODE-re-solve cost.

\end{document}
